\chapter{如何安全改一个实验}

\section{本章要解决的困惑}

读懂系统后，第一反应通常是改参数或加因子。本章只讲安全实验原则，不设计新策略。

\section{一次只改一个问题}

推荐实验模板：

\begin{lstlisting}
Question:
  我想验证什么？

Change:
  只改哪个参数或哪个 gate？

Runtime inputs:
  是否只使用 entry 时刻可见数据？

Profile:
  policy / admission / capacity / exit / leverage

Readout:
  actual_entries / executable return / daily split / worst day
\end{lstlisting}

\section{不要改错地方}

如果只是研究 entry 或 gamma，不要改 Runner。先写 diagnostic 或 fast profile。Runner 是本地交易所，只有在成交、latency、event log 相关时才改。

如果只是看 Monitor，不要让 Monitor 下单。Monitor 只读。

如果只是想引入 Delta/Energy，不要直接说“runtime 已经用了”。先明确它是研究候选，还是 Bot profile 改动。

\section{安全边界}

\begin{lstlisting}
不修改 data/
不修改 date/
不修改 runs/ 中已有用户结果
不打印 .env* 或 rpc.txt
不让 Bot 读 future labels
不让 Runner 读 labels
\end{lstlisting}

\section{实验报告至少写什么}

\begin{longtable}{p{0.26\textwidth}p{0.64\textwidth}}
\toprule
项目 & 内容 \\
\midrule
目标 & 这次验证哪个具体问题 \\
输入 & 使用哪些 runtime-safe 数据 \\
改动 & 改了哪个参数或 profile \\
输出 & actual entries、daily、executable、worst day \\
边界 & fast evidence 还是 strict audit \\
结论 & 是否值得进入下一阶段 \\
\bottomrule
\end{longtable}

\checkpoint{安全实验不是保守不动，而是让每一次结果都能解释、能复现、能审计。}

